\originalchapter{二维晶体与拓扑熔化}
\label{ch:melting}
\source{8.334 L23。}
\originalsection{平移代数序与取向}
二维弹性Gaussian场的位移差方差按$\log r$增长。晶体密度的倒格分量$\rho_G(x)\propto\e^{\ii G\cdot[x-u(x)]}$，故关联$\avg{\rho_G(r)\rho_G(0)^*}\propto r^{-\eta_G}$，$\eta_G$由纵横弹性常数及$G$决定。取向角近$\omega=(\partial_xu_y-\partial_yu_x)/2$，含一个梯度，热长波角涨落在无缺陷晶体中有限。所以二维固体可有代数平移序而长期取向序，并不与位移红外发散矛盾。
\originalsection{位错与位向旋错}
沿闭路位移增量$\oint\partial_i u_j\dd x_i=b_j$为Burgers向量，非零则存在位错；它不能由单值光滑位移场消掉。对六角晶格，键取向绕路改$\pm\pi/3$为基本位向旋错。一个位错可视为相邻反号旋错偶极，两个缺陷层级的能量尺度不同，因而可能两步解束缚。

弹性作用求出的单个位错能$E_d\simeq Yb^2\log(L/a)/(8\pi)+E_c$，$Y=4\mu(\lambda+\mu)/(\lambda+2\mu)$是二维Young模量。位置熵$2\kb\log(L/a)$给长距离重整模量的解束缚条件$Ya^2/(\kb T)=16\pi$，取最短Burgers长度$a$。用裸模量代替已屏蔽模量只是一阶估计。
\originalsection{hexatic与两步图景}
位错解束缚破坏平移代数序，但仍可有六重键角序$\psi_6=\e^{6\ii\theta}$，称hexatic相。其长波角自由能$F_\theta=K_A\int(\nabla\theta)^2/2$。旋错电荷$1/6$，同一涡能熵分析给$K_A/(\kb T)=72/\pi$临界值，$\psi_6$关联指数在临界为$1/4$。再解束缚旋错后为各向同性液体。

KTHNY图景要求缺陷核心代价使连续解束缚流可达相应固定点；某些材料或相互作用可能因核心能低而出现一阶熔化、晶界机制或其他中间结构。因此“所有二维固体必经hexatic”不是该理论的结论。解释实验应同时看平移和取向相关、弹性模量、缺陷密度与有限尺寸，不能仅凭一个结构因子峰宽标成hexatic。
